\section*{Chapter \ref{ch:g4:world-physical-map} --- The Physical Map of the World}

\begin{solution}{exo:g4:world-physical-map:1}
The \omterm{def:g4:world-physical-map:altitude}{altitude} of a place is its height above \omterm{def:g4:world-physical-map:altitude}{sea level}, in metres. It is
measured from \omterm{def:g4:world-physical-map:altitude}{sea level}, the average level of the surface of the \omterm{def:g2:water-on-earth:ocean}{sea}.
\end{solution}

\begin{solution}{exo:g4:world-physical-map:2}
(a) Dark green: between 0 and 200 metres. (b) Pale yellow: between 500 and
1\,000 metres. (c) Almost white: above 4\,000 metres.
\end{solution}

\begin{solution}{exo:g4:world-physical-map:3}
(a) South America. (b) Asia. (c) Africa. (d) North America.
\end{solution}

\begin{solution}{exo:g4:world-physical-map:4}
The highest \omterm{def:g1:landscapes-around-us:landform}{mountain} is Mount Everest (8\,849 metres). The largest \omterm{def:g1:landscapes-around-us:water}{lake} is
the Caspian Sea; the deepest \omterm{def:g1:landscapes-around-us:water}{lake} is \omterm{def:g1:landscapes-around-us:water}{Lake} Baikal.
\end{solution}

\begin{solution}{exo:g4:world-physical-map:5}
The darker the blue, the deeper the \omterm{def:g2:water-on-earth:ocean}{sea}: the \omterm{def:g2:water-on-earth:ocean}{sea} is shallow near the \omterm{def:g1:landscapes-around-us:water}{coasts}
(less than 200 metres) and much deeper in the middle of the \omterm{def:g2:water-on-earth:ocean}{oceans} (more
than 4\,000 metres in many places).
\end{solution}

\begin{solution}{exo:g4:world-physical-map:6}
The \omterm{def:g4:world-physical-map:physical-map}{physical map} shows the \omterm{def:g4:world-physical-map:altitude}{altitude} of the surface. In Antarctica and
Greenland the surface is the top of a thick layer of ice, which lies high
above \omterm{def:g4:world-physical-map:altitude}{sea level}, so they get the colours of high land.
\end{solution}

\begin{solution}{exo:g4:world-physical-map:7}
The bottom of the Challenger Deep (nearly 11\,000 metres below \omterm{def:g4:world-physical-map:altitude}{sea level});
the average floor of the \omterm{def:g2:water-on-earth:ocean}{ocean} (about 3\,700 metres below); a beach (about
0 metres); a village at 1\,500 metres; the top of Mount Everest
(8\,849 metres).
\end{solution}

\begin{solution}{exo:g4:world-physical-map:8}
A \omterm{def:g1:landscapes-around-us:water}{river} is made of many streams that join together, and people do not agree
on which stream is its true beginning. Measured from a farther source, a
\omterm{def:g1:landscapes-around-us:water}{river} is longer. The Nile and the Amazon are so close in length (each about
6\,700 to 6\,800 kilometres) that the answer depends on that choice.
\end{solution}

\begin{solution}{exo:g4:world-physical-map:9}
A quarter of 30 million is seven and a half million square kilometres (30
divided by 4); a third is 10 million (30 divided by 3). The Sahara, eight
and a half million, is more than a quarter of Africa but less than a
third.
\end{solution}

\begin{solution}{exo:g4:world-physical-map:10}
$2500 - 1000 = 1500$: the hut is 1\,500 metres higher. The hut, between
2\,000 and 4\,000 metres, is brown; the \omterm{def:g1:landscapes-around-us:water}{lake}, at 1\,000 metres, is at the
boundary between pale yellow (500 to 1\,000) and light brown (1\,000 to
2\,000): either is accepted.
\end{solution}

\begin{solution}{exo:g4:world-physical-map:11}
10 kilometres is 10\,000 metres. $10000 - 8849 = 1151$: the plane is higher
than Mount Everest, by 1\,151 metres. Over the Challenger Deep, it is
10\,000 metres above the \omterm{def:g2:water-on-earth:ocean}{sea}, and the bottom of the trench is 11\,000
metres below the \omterm{def:g2:water-on-earth:ocean}{sea}: $10000 + 11000 = 21000$ metres above the bottom of
the trench.
\end{solution}
